% Part: computability
% Chapter: computability-theory
% Section: rice-theorem

\documentclass[../../../include/open-logic-section]{subfiles}

\begin{document}

\olfileid{cmp}{thy}{rce}
\olsection{રાઇસનું પ્રમેય}

વિચારો તો દેખાશે કે \olref[tot]{prop:total}ની સાબિતીમાં
$\fn{Tot}$ની કોઈ ખાસ વિશેષતાનો ઉપયોગ થયો નથી. $x$, $K$માં
હોય ત્યારે અને માત્ર ત્યારે $h(x,y)$ અચળ વિધેય $j(y) = 0$
ની જેમ વર્તે એવું આપણે ઘડ્યું; પરંતુ એ સ્થિતિમાં તેને
બીજા કોઈ પણ આંશિક સંગણનીય વિધેયની જેમ વર્તે એવું પણ
ઘડી શકતા. આ અવલોકનથી વધુ સામાન્ય પ્રમેય મળે છે:
આશરે કહીએ તો સંગણનીય વિધેયોનો કોઈ અતુચ્છ ગુણધર્મ
નિર્ણેય નથી.

યાદ રાખો કે $\cfind{0}$, $\cfind{1}$, $\cfind{2}$,~\dots
આંશિક સંગણનીય વિધેયોની આપણી પ્રમાણભૂત પરિગણના છે.

\begin{thm}[રાઇસનું પ્રમેય]
  $C$ આંશિક સંગણનીય વિધેયોનો કોઈ પણ ગણ હોય અને
  $A = \Setabs{n}{\cfind{n} \in C}$ રાખો. જો $A$ સંગણનીય
  હોય, તો $C$ ખાલી છે અથવા $C$ બધા આંશિક સંગણનીય
  વિધેયોનો ગણ છે.
\end{thm}

\emph{સૂચક ગણ} એટલે એવો ગણ $A$ કે જો સૂચકો $n$ અને $m$
એક જ વિધેય ``ગણતા'' હોય, તો બંને $n$ અને $m$ એ $A$માં
હોય, અથવા બંને ન હોય. પ્રમેયમાં આવેલો ગણ~$A$ આ
ગુણધર્મ ધરાવે છે તે જોવું મુશ્કેલ નથી. ઊલટું, જો $A$
સૂચક ગણ હોય અને $C$ આ સૂચકો વડે ગણાતા વિધેયોનો ગણ
હોય, તો $A = \Setabs{n}{\cfind{n} \in C}$.

\begin{explain}
આ પરિભાષામાં રાઇસના પ્રમેયનો અર્થ એવો થાય કે કોઈ અતુચ્છ
સૂચક ગણ નિર્ણેય નથી. પ્રમેય સમજવા \emph{પ્રોગ્રામો}
(માનો તમારી મનપસંદ પ્રોગ્રામિંગ ભાષામાં) અને તેઓ ગણતા
વિધેયો વચ્ચેનો ભેદ ધ્યાનમાં રાખો. પ્રોગ્રામો (સૂચકો)
વાક્યરચનાત્મક પદાર્થો છે; તેમની કેટલીક બાબતો નિશ્ચિતપણે
સંગણનીય છે: આ પ્રોગ્રામમાં ૧૫૦થી વધુ પ્રતીકો છે? તેમાં
૨૨થી વધુ લીટીઓ છે? તેમાં ``while'' વિધાન છે? ``print''
વિધાનના આર્ગ્યુમેન્ટ તરીકે ``hello world'' શબ્દમાળા આવે
છે? રાઇસનું પ્રમેય કહે છે કે પ્રોગ્રામના \emph{વર્તન}
વિશેનો કોઈ અતુચ્છ પ્રશ્ન સંગણનીય નથી. જેમ કે: પ્રોગ્રામ
ઇનપુટ $0$ પર થંભે છે? તે કદી થંભે છે? તે કદી સમ સંખ્યા
આઉટપુટ આપે છે?
\end{explain}

\begin{proof}[રાઇસના પ્રમેયની સાબિતી]
માનો કે $C$ ન તો ખાલી છે, ન બધા આંશિક સંગણનીય વિધેયોનો
ગણ; અને $A$, $C$માંના વિધેયોના સૂચકોનો ગણ છે. $A$
સંગણનીય હોય તો અટકવાની સમસ્યા ઉકેલી શકાય એવું બતાવીશું;
તેથી $A$ સંગણનીય નથી.

સામાન્યતા ગુમાવ્યા વિના, ક્યાંય વ્યાખ્યાયિત ન હોય એવું
વિધેય $f$, $C$માં નથી એમ માની શકીએ (અન્યથા નીચેની
દલીલમાં $C$ અને તેનો પૂરક અદલબદલ કરો). $g$, $C$માંનું
કોઈ પણ વિધેય હોય. વિચાર એવો છે કે $A$નો નિર્ણય કરી
શકીએ, તો $f$ ગણતા સૂચકો અને $g$ ગણતા સૂચકો વચ્ચે
ભેદ કરી શકીએ; પછી આ ક્ષમતાથી અટકવાની સમસ્યા ઉકેલી શકીએ.

રીત આ છે. સર્વવ્યાપી આંશિક સંગણનીય વિધેયોની મદદથી
આવું વિધેય વ્યાખ્યાયિત કરીએ:
\[
h(x,y) \simeq
\begin{cases}
\text{અનિર્ધારિત} & \text{જો $\cfind{x}(x) \fundefined$} \\
g(y) & \text{અન્યથા.}
\end{cases}
\]
$h$ ગણવા પહેલાં $\cfind{x}(x)$ ગણવાનો પ્રયત્ન કરીએ.
તે સંગણના થંભે, તો પછી~$g(y)$ ગણીએ; અને \emph{તે}
સંગણના પણ થંભે, તો આઉટપુટ પાછું આપીએ. વધુ ઔપચારિક
રીતે આમ લખી શકાય:
\[
h(x,y) \simeq \Proj{2}{0}(g(y),\fn{Un}(x,x)).
\]
અહીં $\Proj{2}{0}(z_0, z_1) = z_0$ એ $2$-સ્થાનીય પ્રક્ષેપ
વિધેય છે, જે $0$મા આર્ગ્યુમેન્ટ આપે છે અને
સંગણનીય છે.

તેથી $h$ આંશિક સંગણનીય વિધેયોના સંયોજનથી બને છે અને
જમણી બાજુ ત્યારે અને માત્ર ત્યારે વ્યાખ્યાયિત થઈને~$g(y)$
જેટલી થાય છે જ્યારે $\fn{Un}(x,x)$ અને $g(y)$ બંને
વ્યાખ્યાયિત હોય.

નોંધો કે કોઈ નિશ્ચિત~$x$ માટે $\cfind{x}(x)$ અનિર્ધારિત
હોય, તો દરેક~$y$ માટે $h(x,y)$ અનિર્ધારિત છે; અને
$\cfind{x}(x)$ વ્યાખ્યાયિત હોય, તો $h(x,y) \simeq g(y)$.
આમ~$x$નું કોઈ પણ નિશ્ચિત મૂલ્ય લેતાં $h(x,y)$ કાં તો
$f$ની જેમ અથવા $g$ની જેમ વર્તે છે. એટલે
$\cfind{x}(x)$ વ્યાખ્યાયિત છે કે નહીં તે નક્કી કરવું
આ બે સ્થિતિમાંથી કઈ છે તે નક્કી કરવા સમાન છે. પરંતુ
તે $h_x(y) \simeq h(x,y)$ એ~$C$માં છે કે નહીં તે
નક્કી કરવા સમાન છે; $A$ સંગણનીય હોય તો તે કરી શકીએ.

વધુ ઔપચારિક રીતે, $h$ આંશિક સંગણનીય હોવાથી કોઈ
સૂચક~$e$ માટે તે વિધેય $\cfind{e}[2]$ છે.
\sourcecorrection{OLCT-010}{સ્થિર અંગ્રેજી મૂળ અહીં
દ્વિ-આર્ગ્યુમેન્ટવાળા $h(x,y)$ માટે એક-આર્ગ્યુમેન્ટવાળું
$\cfind{e}$ લખે છે. એ જ સૂચકનું દ્વિ-આર્ગ્યુમેન્ટવાળું
સ્વરૂપ $\cfind{e}[2]$ અહીં વપરાયું છે.}
$s$-$m$-$n$ પ્રમેયથી એવું આદિમ પુનરાવર્તી વિધેય~$s$
છે કે દરેક~$x$ માટે $\cfind{s(e,x)}(y) = h_x(y)$.
હવે દરેક $x$ માટે $\cfind{x}(x) \fdefined$ હોય, તો
$\cfind{s(e,x)}$ એ~$g$ જેવું જ વિધેય છે અને તેથી
$s(e,x)$ એ $A$માં છે. બીજી તરફ, $\cfind{x}(x)
\uparrow$ હોય, તો $\cfind{s(e,x)}$ એ~$f$ જેવું જ
વિધેય છે અને તેથી $s(e,x)$ એ~$A$માં નથી. બીજા
શબ્દોમાં, દરેક~$x$ માટે $x \in K$ ત્યારે અને માત્ર
ત્યારે જ્યારે $s(e,x) \in A$. $A$ સંગણનીય હોત,
તો $K$ પણ હોત; આ વિરોધાભાસ છે. તેથી $A$
સંગણનીય નથી.
\end{proof}

રાઇસનું પ્રમેય ખૂબ શક્તિશાળી છે. આગળનું તાત્કાલિક પરિણામ
તેના કેટલાક ઉપયોગો બતાવે છે.
\begin{cor}
આગળના ગણો અનિર્ણેય છે.
\begin{enumerate}
\item $\Setabs{x}{\text{$17$ એ $\cfind{x}$ના મૂલ્યગણમાં છે}}$
\item $\Setabs{x}{\text{$\cfind{x}$ અચળ છે}}$
\item $\Setabs{x}{\text{$\cfind{x}$ સંપૂર્ણ છે}}$
\item $\Setabs{x}{\text{જ્યારે $y < y'$, ત્યારે $\cfind{x}(y) \fdefined$,
    અને $\cfind{x}(y') \fdefined$ હોય તો $\cfind{x}(y) < \cfind{x}(y')$}}$
\end{enumerate}
\end{cor}

\begin{proof}
  આ બધા અતુચ્છ સૂચક ગણો છે.
\end{proof}

\end{document}
